\section*{Chapitre \ref{ch:g12:arith} --- Arithmétique}

\begin{solution}{exo:g12:arith:1}
$17 \times 119 = 2023$, donc $2026 = 17 \times 119 + 3$~: quotient $119$,
reste $3$. Pour $-2026$~: $-2026 = 17\times(-120) + 14$ (en effet
$17 \times 120 = 2040$ et $2040 - 2026 = 14$)~: quotient $-120$, reste
$14$ (le reste doit être dans $\intco{0}{17}$, donc ce n'est \emph{pas}
$-3$).
\end{solution}

\begin{solution}{exo:g12:arith:2}
Modulo $10$~: $7^2 = 49 \equiv 9 \equiv -1$. D'où
$7^{100} = \left(7^2\right)^{50} \equiv (-1)^{50} = 1 \pmod{10}$~: le
dernier chiffre de $7^{100}$ est $1$.
\end{solution}

\begin{solution}{exo:g12:arith:3}
Euclide~: $1071 = 2\times462 + 147$~; $462 = 3\times147 + 21$~;
$147 = 7\times21 + 0$. Donc $\gcd = 21$.

Remontée~: $21 = 462 - 3\times147 = 462 - 3(1071 - 2\times462)
= 7\times462 - 3\times1071$. Ainsi $u = -3$, $v = 7$~:
$1071\times(-3) + 462\times7 = 21$.
\end{solution}

\begin{solution}{exo:g12:arith:4}
Tout entier est $\equiv 0, 1, 2$ ou $3 \pmod 4$, et en élevant au carré~:
$0^2 \equiv 0$, $1^2 \equiv 1$, $2^2 = 4 \equiv 0$, $3^2 = 9 \equiv 1$.
Donc $n^2 \equiv 0$ ou $1 \pmod 4$. Une somme de deux carrés est alors
congrue à $0 + 0$, $0 + 1$ ou $1 + 1$, \emph{c.-à-d.} à $0$, $1$ ou
$2 \pmod 4$ --- jamais à $3$.
\end{solution}

\begin{solution}{exo:g12:arith:5}
\omterm{def:g12:arith:divides}{Divisibilité} par $2$~: parmi $n$ et $n + 1$, l'un est pair. \omterm{def:g12:arith:divides}{Divisibilité}
par $3$~: si $n \equiv 0$, alors $3 \mid n$~; si $n \equiv 1 \pmod 3$,
alors $2n + 1 \equiv 3 \equiv 0$~; si $n \equiv 2$, alors
$n + 1 \equiv 0$. Dans tous les cas $3$ divise le produit. Comme
$\gcd(2,3) = 1$, le \cref{cor:g12:arith:coprimeprod} donne
$6 \mid n(n+1)(2n+1)$.
(Cela reprouve aussi que $\frac{n(n+1)(2n+1)}{6}$, la somme des carrés de
l'\cref{exo:g12:seq:1}, est un entier.)
\end{solution}

\begin{solution}{exo:g12:arith:6}
On cherche l'inverse de $5$ modulo $11$~: en testant (ou Bézout),
$5 \times 9 = 45 = 44 + 1 \equiv 1 \pmod{11}$. En multipliant la
\omterm{def:g12:arith:congruence}{congruence} par $9$~:
\[
x \equiv 9 \times 3 = 27 \equiv 5 \pmod{11}.
\]
Les solutions sont les entiers $x = 5 + 11k$, $k \in \Z$. (Vérification~:
$5\times5 = 25 \equiv 3 \pmod{11}$.)
\end{solution}

\begin{solution}{exo:g12:arith:7}
$\gcd(17, 40) = 1$, donc des solutions existent. Euclide~:
$40 = 2\times17 + 6$~; $17 = 2\times6 + 5$~; $6 = 5 + 1$. En remontant~:
$1 = 6 - 5 = 6 - (17 - 2\times6) = 3\times6 - 17
= 3(40 - 2\times17) - 17 = 3\times40 - 7\times17$. D'où
$17\times(-7) - 40\times(-3) = 1$~: la solution particulière
$(x_0, y_0) = (-7, -3)$.

Solution générale de $17x - 40y = 1$~: en soustrayant la relation
particulière, $17(x + 7) = 40(y + 3)$~; comme $\gcd(17, 40) = 1$, Gauss
donne $40 \mid x + 7$, donc $x = -7 + 40k$ puis $y = -3 + 17k$,
$k \in \Z$ (qui toutes vérifient).

Pour $17x - 40y = 6$, multiplier la solution particulière par $6$~:
$(x_1, y_1) = (-42, -18)$, et le même raisonnement donne
\[
x = -42 + 40k, \qquad y = -18 + 17k, \qquad k \in \Z .
\]
(Par ex.\ $k = 2$~: $x = 38$, $y = 16$~; en effet $17\times38 - 40\times16
= 646 - 640 = 6$.)
\end{solution}

\begin{solution}{exo:g12:arith:8}
Supposer $\sqrt2 = \frac ab$ avec $a, b \in \N^*$~; alors $a^2 = 2b^2$.
Dans la factorisation d'un carré, tout exposant est pair~; donc
l'exposant de $2$ dans $a^2$ est pair, tandis que dans $2b^2$ il est
impair (un de plus qu'un nombre pair). Deux factorisations du même entier
avec des exposants de $2$ différents contredisent l'unicité du
\cref{thm:g12:arith:fta}. Donc aucune telle fraction n'existe~:
$\sqrt2 \notin \Q$.
\end{solution}

\begin{solution}{exo:g12:arith:9}
\emph{1.} De $k\binom pk = p \binom{p-1}{k-1}$, $p$ divise
$k\binom pk$. Pour $1 \leq k \leq p-1$, $p \nmid k$ et $p$ \omterm{def:g12:arith:prime}{premier} donnent
$\gcd(p, k) = 1$, donc le lemme de Gauss fournit $p \mid \binom pk$.

\emph{2.} Récurrence sur $a$. Pour $a = 0$~: $0^p \equiv 0$. Supposer
$a^p \equiv a \pmod p$. Par la formule du binôme,
\[
(a+1)^p = \sum_{k=0}^{p} \binom pk a^k
\equiv a^p + 1 \pmod p,
\]
tous les termes du milieu s'annulant modulo $p$ par le point~1. Par
l'hypothèse de récurrence, $(a+1)^p \equiv a + 1 \pmod p$. Cela prouve
$a^p \equiv a$ pour tout $a \in \N$, et le cas $a < 0$ suit en écrivant
$a \equiv a + kp$ pour un représentant positif convenable.
\end{solution}

\begin{solution}{exo:g12:arith:10}
$n \equiv 2 \pmod 3$ et $n \equiv 3 \pmod 5$~: écrire $n = 2 + 3s$~; alors
$2 + 3s \equiv 3 \pmod 5$, \emph{c.-à-d.} $3s \equiv 1 \pmod 5$.
L'inverse de $3$ modulo $5$ est $2$ ($3\times2 = 6 \equiv 1$), donc
$s \equiv 2 \pmod 5$, disons $s = 2 + 5t$, et $n = 8 + 15t$~: les deux
premières conditions signifient $n \equiv 8 \pmod{15}$.

En ajoutant $n \equiv 2 \pmod 7$~: $8 + 15t \equiv 2 \pmod 7$, et
$15 \equiv 1 \pmod 7$, donc $t \equiv -6 \equiv 1 \pmod 7$, disons
$t = 1 + 7u$. D'où $n = 23 + 105u$~:
\[
n \equiv 23 \pmod{105}.
\]
(Vérification~: $23 = 3\times7 + 2 = 5\times4 + 3 = 7\times3 + 2$.)
\end{solution}

\begin{solution}{pb:g12:arith:1}
\textbf{1.} $2026 = 289 \times 7 + 3$, donc $2026 \equiv 3
\pmod 7$. Puissances de $7$ modulo $10$~: $7, 9, 3, 1$, cycle de
longueur $4$~; comme $100 \equiv 0 \pmod 4$, le chiffre des unités
de $7^{100}$ est $1$.

\textbf{2.} $5^2 = 25 \equiv -1 \pmod{13}$, donc
$5^{116} = \left(5^2\right)^{58} \equiv (-1)^{58} = 1$ et
$5^{117} \equiv 5 \pmod{13}$.

\textbf{3.} L'inverse de $3$ modulo $7$ est $5$ (car
$15 \equiv 1$)~: $x \equiv 5 \times 5 = 25 \equiv 4 \pmod 7$.

\textbf{4.} $97 = 2 \times 35 + 27$~; $35 = 27 + 8$~;
$27 = 3 \times 8 + 3$~; $8 = 2 \times 3 + 2$~; $3 = 2 + 1$. En
remontant~: $1 = 97 \times 13 + 35 \times (-36)$. Donc
$35 \times (-36) \equiv 1 \pmod{97}$~: l'inverse de $35$ est
$-36 \equiv 61 \pmod{97}$.

\textbf{5.} $a$ est inversible modulo $n$ exactement lorsque
$\gcd(a, n) = 1$~: Bézout fournit $au + nv = 1$, c'est-à-dire
$au \equiv 1$~; réciproquement, l'existence d'un inverse force le
\omterm{def:g12:arith:gcd}{pgcd} à diviser $1$.

\textbf{6.} $0{\cdot}10 + 3{\cdot}9 + 0{\cdot}8 + 6{\cdot}7 +
4{\cdot}6 + 0{\cdot}5 + 6{\cdot}4 + 1{\cdot}3 + 5{\cdot}2 +
2{\cdot}1 = 132 = 12 \times 11 \equiv 0 \pmod{11}$~: le code est
valide.

\textbf{7.} La somme change de $wd$ avec $1 \leq w \leq 10$ et
$1 \leq \abs d \leq 9$~: comme $11$ est \omterm{def:g12:arith:prime}{premier} et ne divise aucun
des deux facteurs, il ne peut pas diviser le produit
(\cref{thm:g12:arith:gauss} et
\cref{prop:g12:arith:primedivides})~: la somme modifiée n'est
jamais de nouveau $\equiv 0$, et toute erreur sur un seul chiffre
déclenche l'alarme.

\textbf{8.} Échanger deux chiffres adjacents $a$ et $b$ (de poids
$w + 1$ et $w$) change la somme de
$(w+1)b + wa - (w+1)a - wb = b - a \not\equiv 0$ dès que
$a \neq b$~: c'est détecté. Le secret est la \emph{primalité} de
$11$~: modulo $10$, des produits comme $5 \times 2$ s'annulent sans
qu'aucun facteur ne soit nul, si bien qu'une erreur de $\pm 2$ sur
un poids $5$ (ou une transposition malchanceuse) pourrait passer.

\textbf{9.} Somme pondérée des douze chiffres~: $119$~; la clé
doit la compléter jusqu'à un multiple de $10$, c'est donc $1$ (code
complet $978\,2940199\,051$). L'EAN rate les transpositions
adjacentes vérifiant $2(a - b) \equiv 0 \pmod{10}$, c'est-à-dire
$\abs{a - b} = 5$~: échanger un $2$ et un $7$, par exemple, passe
inaperçu --- c'est le prix du sympathique \omterm{def:g12:complex:modulus}{module} $10$.

\textbf{10.} En doublant un chiffre sur deux à partir de la droite
et en repliant ($16 \to 7$, etc.), la somme vaut $80 \equiv 0
\pmod{10}$~: la carte de test est validée.

\textbf{11.} Avec un \omterm{def:g12:complex:modulus}{module} \omterm{def:g12:arith:prime}{premier}, tous les poids sont
inversibles, si bien que \emph{toutes} les erreurs simples et
\emph{toutes} les transpositions adjacentes sont attrapées ---
c'est le luxe de l'ISBN~; les dispositifs modulo $10$ conservent
des chiffres commodes pour l'humain et acceptent un petit angle
mort.

\textbf{12.} $2^{10} = 1024 = 3 \times 341 + 1 \equiv 1
\pmod{341}$, donc $2^{340} = \left(2^{10}\right)^{34} \equiv 1$. Et
pourtant $341 = 11 \times 31$ est composé~: il passe le test de
Fermat en base $2$ sans être \omterm{def:g12:arith:prime}{premier}. Morale~: la \omterm{def:g12:arith:congruence}{congruence} de
Fermat est nécessaire, non suffisante --- tester la primalité
réclame des outils plus fins (et en reçoit, dans les volumes
universitaires).

\textbf{13.} $3d \equiv 1 \pmod{20}$~: $d = 7$ (car
$21 = 20 + 1$).

\textbf{14.} $c = 4^3 = 64 \equiv 31 \pmod{33}$.

\textbf{15.} $31 \equiv -2$~: $(-2)^7 = -128$, et
$-128 + 4 \times 33 = 4$~: le chiffré se déchiffre en $m = 4$. La
serrure tourne.

\textbf{16.} Modulo $3$~: si $3 \nmid m$, alors $m^2 \equiv 1$
(Fermat), donc $m^{21} = m \cdot \left(m^2\right)^{10} \equiv m$~;
si $3 \mid m$, les deux membres sont $\equiv 0$. Modulo $11$~:
$m^{10} \equiv 1$ ou bien $11 \mid m$, et
$m^{21} = m \cdot \left(m^{10}\right)^2 \equiv m$. Ainsi $3$ et
$11$ divisent tous deux $m^{21} - m$, et comme ils sont \omterm{def:g12:arith:gcd}{premiers
entre eux}, leur produit $33$ le divise aussi (Gauss)~:
$m^{21} \equiv m \pmod{33}$. L'exposant $ed = 21 = 1 + 20k$ a été
bâti pour que les deux exposants de Fermat ($2$ et $10$, qui
divisent $20$) disparaissent.

\textbf{17.} Multiplier deux \omterm{def:g12:arith:prime}{nombres premiers} de $300$ chiffres
prend une microseconde~; les retrouver à partir de leur produit met
en échec tout algorithme connu et tous les ordinateurs du monde ---
la serrure est une rue à sens unique. (Notre $n = 33$ est cette rue
à l'échelle du jouet, parcourable dans les deux sens.)

\textbf{18.} En testant les restes (ou en construisant par
Bézout)~: $n \equiv 8 \pmod{15}$, d'où les effectifs
$8, 23, 38, 53, \dots$ Unicité modulo $15$~: deux solutions
diffèrent d'un multiple de $3$ et d'un multiple de $5$, donc de
$15$ ($3$ et $5$ étant \omterm{def:g12:arith:gcd}{premiers entre eux}, Gauss). Le général qui
commande à $1000$ soldats annonce «~$8$~» après trois
rassemblements rapides --- l'antique astuce de comptage.

\textbf{19.} $10 \equiv 1 \pmod 9$ donne $10^k \equiv 1$, donc
$\sum d_k 10^k \equiv \sum d_k$~: un nombre et la somme de ses
chiffres sont congrus modulo $9$ (et modulo $3$). Et
$10 \equiv -1 \pmod{11}$ donne
$\sum d_k 10^k \equiv \sum (-1)^k d_k$~: c'est la règle alternée.
Deux règles d'enfance, une ligne chacune.

\textbf{20.} Les \omterm{def:g12:arith:congruence}{congruences} ont fait des restes une arithmétique
(partie I)~; Bézout a frappé les inverses qui résolvent les
\omterm{def:g12:arith:congruence}{congruences} linéaires et fournissent le $d$ de RSA (questions 4 et
13)~; le petit théorème de Fermat a ouvert et refermé la serrure
(questions 15 et 16)~; le recollement de \omterm{def:g12:complex:modulus}{modules} \omterm{def:g12:arith:gcd}{premiers entre eux}
a compté les soldats et achevé la démonstration (questions 16 et
18). Verdict sur Hardy~: le plus pur théorème qu'il connût garde
aujourd'hui chacun de nos achats --- la pureté, avec le temps, est
la chose la plus applicable qui soit.
\end{solution}
